\documentclass[11pt,a4paper]{article}

\usepackage[T1]{fontenc}
\usepackage[utf8]{inputenc}
\usepackage{amsmath,amssymb,amsthm}
\usepackage{mathtools}
\usepackage{hyperref}
\usepackage{doi}
\usepackage{geometry}
\usepackage{booktabs}
\usepackage{array}
\usepackage{enumitem}

\geometry{margin=2.5cm}

\hypersetup{
  colorlinks=true,
  linkcolor=black,
  citecolor=black,
  urlcolor=black
}

\theoremstyle{definition}
\newtheorem{remark}{Remark}

\newcommand{\Heis}{\mathrm{Heis}}
\newcommand{\SU}{\mathrm{SU}}
\newcommand{\su}{\mathfrak{su}}
\newcommand{\heis}{\mathfrak{heis}}
\newcommand{\BI}{\mathrm{BI}}
\newcommand{\BFS}{\mathrm{BFS}}
\newcommand{\Leff}{L_{\mathrm{eff}}}
\newcommand{\Lpi}{L_\Pi}
\newcommand{\geff}{g^{\mu\nu}}
\newcommand{\Atau}{A_\tau}
\newcommand{\AH}{A_H}
\newcommand{\AZ}{A_Z}
\newcommand{\Dhom}{D_{\mathrm{hom}}}
\newcommand{\Sym}{\mathrm{Sym}}
\newcommand{\ImH}{\operatorname{Im}\mathbb{H}}
\renewcommand{\Im}{\operatorname{Im}}

\title{The Emergent Geometry Sub-Programme\\[0.4em]
\large Presentation Note 2}

\author{J\'{e}r\^{o}me Beau\\[0.3em]
\small Independent Researcher\\
\small \texttt{jerome.beau@cosmochrony.org}}

\date{2026}

\begin{document}

  \maketitle

  \begin{abstract}
    The emergent geometry sub-programme of the Cosmochrony corpus addresses a single
    central question: how does the admissible Weil--Heisenberg fibre give rise to an effective
    four-dimensional Lorentzian geometry?
    Starting from the admissibility filter $\Pi_q$ acting on the Weil representation
    $V_\rho$ of $\Heis_3(\mathbb{Z}/q\mathbb{Z})$, the sub-programme closes the following chain:
    \[
      \Pi_q
      \;\Longrightarrow\;
      V_\rho \simeq L^2(\mathbb{Z}/q\mathbb{Z})
      \;\Longrightarrow\;
      \Heis_3(\mathbb{R})
      \;\Longrightarrow\;
      \Leff
      \;\Longrightarrow\;
      \geff
      \;\Longrightarrow\;
      \geff = 2\eta^{\mu\nu}.
    \]
    This note maps eleven constituent papers (Q5a, Q5a-O2, Q5b, Q6b, Q7--Q11, U1, W1, H2)
    across five internal phases and records the status of every result.
    The remaining open items are not obstructions to the admissible-sector metric closure.
    They concern, respectively, the finite-$q$ bridge formulation and the extension of the
    Q5a Mosco tightness control from the admissible sector to the full function space.
  \end{abstract}

  \medskip
  \noindent\textbf{Keywords.}
  emergent spacetime; Heisenberg group; Carnot geometry; Mosco convergence; Lorentzian
  signature; effective metric; sub-Riemannian geometry; Casimir rigidity; non-injective
  projection; admissibility; presentation note.

  \newpage
  \tableofcontents

  \newpage

  \section{Motivation and Central Question}
    \label{sec:motivation}

    The emergent geometry sub-programme answers the following question.

    \medskip
    \noindent\textit{Central question.}
    The admissibility filter $\Pi_q$ of the Cosmochrony framework acts on the Weil
    representation $V_\rho \simeq L^2(\mathbb{Z}/q\mathbb{Z})$ of
    $\Heis_3(\mathbb{Z}/q\mathbb{Z})$, selecting admissible modes according to the
    Born--Infeld bounded-flux constraint.
    \textit{What effective geometry is selected in the large-$q$ limit of this structure,
      and what determines the numerical values of its metric coefficients?}

    \medskip
    The practical importance of this question is that the entire gravitational, gauge, and
    fermionic content of Branch~III depends on the identification of an effective spacetime.
    Without a derived metric, the dynamical outputs of the Gravity paper~\cite{Beau2026h},
    the gauge structure of Q6a~\cite{Beau2026q6a}, and the spinorial structure of
    Q14~\cite{Beau2026q14} would each require an independently postulated background.
    The emergent geometry sub-programme removes this postulate: spacetime is a
    forced consequence of admissibility, not an input.

    \medskip
    This note is concerned exclusively with the reconstruction of the effective metric
    $\geff = 2\eta^{\mu\nu}$.
    The dynamical content — why the effective metric satisfies the Einstein equations — is
    the subject of the spectral gravity sub-programme and is not claimed here.

  \section{Position in the Cosmochrony Programme}
    \label{sec:position}

    The emergent geometry sub-programme occupies the interface between Branch~I (axiomatic
    primitive) and Branch~III (physical observables).
    It takes the algebraic output of Branch~I — the Weil representation of
    $\Heis_3(\mathbb{Z}/q\mathbb{Z})$ and the Born--Infeld admissibility constraint — and
    produces the effective Lorentzian metric used throughout Branch~III.

    \medskip
    Its relation to the spectral admissibility sub-programme (Presentation Note~1) is the
    following.
    Note~1 establishes that the admissible sector is the spin-$\tfrac{1}{2}$ sector
    $V_\rho \cong \mathbb{C}^2$ and that $\Sigma_c(n_3) = 3$ with $\Im\mathbb{H} \cong \su(2)$.
    The emergent geometry sub-programme takes these facts as inputs and asks: what metric
    structure do they support?

    \begin{remark}[Vocabulary note]
Throughout this note, the term \emph{admissibility constraint} replaces the earlier
``relaxation'' vocabulary of the preliminary Cosmochrony formulation.
The substrate $\chi$ is static: no substrate dynamics or relaxation mechanism is
postulated at the foundational level.
Observable evolution is an induced ordering of admissible projections, not a dynamics of
$\chi$ itself.
    \end{remark}

    \begin{remark}[Scope: metric reconstruction only]
This sub-programme closes the chain $\Pi_q \to \geff = 2\eta^{\mu\nu}$.
It does not address the dynamics of $\geff$: the question of why the effective metric
satisfies the Einstein equations is handled by the spectral gravity sub-programme, which
takes the metric established here as its starting point.
    \end{remark}

  \section{Logical Chain of the Sub-Programme}
    \label{sec:chain}

    The sub-programme is organised around the following derivation chain.

    \begin{equation}
      \label{eq:chain}
      \underbrace{\Pi_q}_{\text{admissibility filter}}
      \;\Longrightarrow\;
      \underbrace{\Lpi = -A\partial_x^2}_{\substack{\text{1D shadow} \\ \text{(Mosco, Q5a)}}}
      \;\Longrightarrow\;
      \underbrace{\Leff \text{ on } \mathbb{R}_\tau \times \Heis_3(\mathbb{R})}_{\substack{\text{4D effective operator} \\ \text{(Carnot, Q5b)}}}
      \;\Longrightarrow\;
      \underbrace{\geff \propto A_{\mu\nu}}_{\substack{\text{metric extraction} \\ \text{(Q5b + Q9)}}}
      \;\Longrightarrow\;
      \underbrace{\geff = 2\eta^{\mu\nu}}_{\substack{\text{metric closure} \\ \text{(Q8--Q11, with U1, W1, H2)}}}.
    \end{equation}

    The chain has five conceptually distinct stages, each resolved by a distinct group of papers.

    \begin{enumerate}
      \item \textit{Discrete-to-continuum.}
      Q5a establishes the Mosco convergence of the admissibility Dirichlet forms $\mathcal{E}_q$
      on $\mathrm{Cay}(G_q, S_q)$ to a one-dimensional shadow operator
      $\Lpi = -A\partial_x^2$ on $L^2(\mathbb{R})$.
      Q5a-O2 closes the two remaining open hypotheses of Q5a ([H-E1] and Conjecture~[C])
      via spectral atomicity.
      H2 closes hypothesis~[H2].

      \item \textit{Dimensional promotion.}
      Q5b promotes the one-dimensional $\Lpi$ to a four-dimensional effective operator
      $\Leff$ on $\mathbb{R}_\tau \times \Heis_3(\mathbb{R})$ via the Bass--Guivarc'h
      homogeneous dimension $\Dhom = 4$ and the Carnot convergence of BFS shells.

      \item \textit{Metric extraction.}
      Q5b Theorem~5.2, made unconditional by Q9's proof of [H-lift], extracts the effective
      co-metric tensor $\geff \propto A_{\mu\nu}$ from the principal symbol of $\Leff$.
      Q5b Theorem~6.1 derives Lorentzian signature $(-, +, +, +)$ from the Born--Infeld
      admissibility constraint.

      \item \textit{Coefficient determination.}
      Q7 identifies the effective projection space $H_{\mathrm{eff}} = \Sym^2(V_\rho)$ as
      the unique irreducible three-dimensional $\su(2)$-module and establishes the bridge
      rigidity $\phi: \Sym^2(V_\rho) \xrightarrow{\sim} W_{\mathrm{sp}}$.
      Q8 derives $\AZ = C_{\su(2)} = 2$ via Casimir rigidity on the central direction.
      Q10 establishes $\AH \to 2$ under spectral universality~[U].
      U1 proves [U] with rate $\varepsilon(q) = O(q^{-1/2})$.

      \item \textit{Metric closure.}
      Q11 closes the sole remaining open coefficient $\Atau = 2$ via temporal Casimir
      rigidity, giving $\geff = 2\eta^{\mu\nu}$.
      W1 closes [H-w].
    \end{enumerate}

  \section{Constituent Papers}
    \label{sec:papers}

    \subsection{Continuum limit and one-dimensional shadow: Q5a and Q5a-O2}
      \label{sec:q5a}

      \textbf{Q5a}~\cite{Beau2026q5a} establishes, conditionally on hypotheses [H1], [H-w],
      [H-E1], and Conjecture~[C], three convergences.
      (T1)~A Hilbert inductive limit $\varinjlim_q \mathcal{C}_q \simeq L^2(\mathbb{R})$, where
      $\mathcal{C}_q \simeq L^2(\mathbb{Z}/q\mathbb{Z})$ is the Weil representation space at
      finite $q$.
      (T2)~Convergence of rescaled Weil generators to the metaplectic generators of
      $\Heis_3(\mathbb{R})$.
      (T3)~Mosco convergence of the admissibility Dirichlet forms $\mathcal{E}_q$ to the
      second-order operator $\Lpi = -A\partial_x^2$ on $L^2(\mathbb{R})$, $A > 0$.

      The key conceptual content of Q5a is that the continuum limit is not an additional
      assumption: it is forced by the coercivity and compactness properties of the admissible
      forms.
      Uniform energy bounds prevent spectral mass from escaping to high frequencies;
      the Mosco limit is the unique operator compatible with spectral stability.

      $\Lpi$ is a one-dimensional operator.
      Q5a states explicitly that ``the operator $\Lpi$ is the one-dimensional operator-theoretic
      shadow of the homogeneous four-dimensional structure of $\Heis_3(\mathbb{R})$.
      The identification of $\Lpi$ with a four-dimensional effective geometry is the subject of
      Q5b.''

      The O-series link is structural: Hypothesis~[H-E1] (uniform spectral gap) is equivalent to the
      stabilisation of $\delta_{\mathrm{pair}}$ across primes, supported by the O12--O25 campaign.
      Hypothesis~[H-w] (admissibility weight convergence $a_q(s) \to A > 0$) is proved in
      W1~\cite{Beau2026w1}.

      \textbf{Q5a-O2}~\cite{Beau2026q5ao2} resolves the two remaining open hypotheses of Q5a.
      It proves that the admissible sector $\mathrm{ran}(\Pi_q)$ is spectrally atomic:
      it is spanned by three pure Fourier modes per conjugate pair, with macroscopic frequencies
      $\xi_c^{(q)}/q \to \omega_c \neq 0$.
      From this structure, coercivity at the optimal diffusive scale follows without any Nash or
      Rellich inequality, and Mosco compactness follows from Bolzano--Weierstrass in $\mathbb{C}^2$.
      Hypotheses [H-E1] and [C] become unconditional theorems on the admissible sector.

    \subsection{Carnot geometry and Lorentzian signature: Q5b}
      \label{sec:q5b}

      \textbf{Q5b}~\cite{Beau2026q5b} delivers three results that together promote the
      one-dimensional shadow $\Lpi$ to a four-dimensional effective geometry.

      \textit{Carnot convergence} (Theorem~3.2; structural).
      The BFS shell stratification $\{S_n\}$ of $G_q = \Heis_3(\mathbb{Z}/q\mathbb{Z})$
      converges, in the pre-saturation regime $n \ll q$, to the Carnot--Carath\'{e}odory
      sphere foliation of $\Heis_3(\mathbb{R})$.
      The homogeneous dimension $\Dhom = 4$ (Bass--Guivarc'h) gives the limiting geometry
      the spectral and volume-growth properties of a four-dimensional space.
      This result is structural: it follows from Pansu's asymptotic cone theorem and requires
      no new hypothesis.

      \textit{Metric extraction} (Theorem~5.2; unconditional via Q9~\cite{Beau2026q9}).
      Under [H-lift] — proved in Q9 as a theorem rather than a hypothesis — $\Lpi$ is the image,
      under the Schr\"{o}dinger representation $\pi_1$, of the kinetic sector of the
      sub-Riemannian Laplacian $\Delta_H = \tilde{X}^2 + \tilde{Y}^2$ on $\Heis_3(\mathbb{R})$,
      restricted to the admissible sector.
      The effective co-metric $\geff \propto A_{\mu\nu}$ is extracted from the principal symbol
      of the resulting effective operator $\Leff$ on $\mathbb{R}_\tau \times \Heis_3(\mathbb{R})$.

      \textit{Lorentzian signature} (Theorem~6.1; conditional on Theorem~5.2).
      The Born--Infeld admissibility constraint forces signature $(-, +, +, +)$ by exclusion:
      Riemannian signature is excluded by the finite-capacity bound of A4; ultra-hyperbolic
      signature fails to admit a globally well-posed initial value problem; degenerate
      signature is excluded by sub-principal completion; parabolic signature implies a
      characteristic direction incompatible with monotonic projective ordering.

      After normalising the maximal admissible propagation speed to $c = 1$, the effective
      metric is fixed to $\geff = \mathrm{diag}(-1, +1, +1, +1)$ in the homogeneous isotropic
      regime.

    \subsection{Effective spacetime output: Q6b}
      \label{sec:q6b}

      \textbf{Q6b}~\cite{Beau2026q6b} is the integrative output of the sub-programme.
      Its role is not to re-derive the geometric convergence theorems of Q5a--Q5b, but to assemble
      the chain into the macroscopic effective-spacetime statement:
      \[
        \Pi_q \;\longrightarrow\; \Leff \;\longrightarrow\; \geff \;\longrightarrow\; G_{\mu\nu}.
      \]

      Three results are obtained.
      (i)~The Schwarzschild metric as the unique stationary, spherically symmetric exterior
      solution compatible with the admissibility flow, obtained from flux conservation within the
      effective geometry.
      (ii)~A structural interpretation of the Schwarzschild horizon as a degeneracy of the
      principal symbol of $\Leff$, not a singularity of the underlying admissible structure.
      (iii)~The identification of the Einstein equations as consistency conditions of the emergent
      geometry, via the spectral entropy functional of the Gravity paper.

      Q6b closes the chain $\Pi_q \to \Leff \to \geff \to G_{\mu\nu}$ under the Q5a--Q5b
      hypotheses.
      With [H-lift] proved in Q9, the results of Q6b are unconditional with respect to [H-lift].
      The remaining conditionality stems from the Q5a Mosco hypotheses [H1], [H-w], [H-E1],
      and~[C]; of these, [H-w] is proved in W1, and [H-E1] and [C] are closed in Q5a-O2.

      \begin{remark}[Status of Q6b]
Q6b is listed as published in the programme.
Its status as \emph{integrative output} means that it assembles rather than proves: it
inherits the metric from Q5b and passes the Einstein equations to the Gravity paper.
It should not be cited as the primary source for either the metric or the Einstein
equations; those primary sources are Q5b and the Gravity paper respectively.
      \end{remark}

    \subsection{Spatial, vertical, horizontal, and temporal coefficients: Q7--Q11}
      \label{sec:q7q11}

      Papers Q7--Q11 determine the four numerical entries of the effective metric by distinct
      representation-theoretic arguments, all converging to the same value $A = 2$.

      \textbf{Q7}~\cite{Beau2026q7} establishes three structural results.
      First, $\su(2) \not\cong \heis_3(\mathbb{R})$ as Lie algebras, so the bridge
      $\phi: H_{\mathrm{eff}} \to W_{\mathrm{sp}}$ operates at the representation level.
      Second, the effective projection space $H_{\mathrm{eff}} = \Sym^2(V_\rho)$ carries
      the unique irreducible three-dimensional $\su(2)$-module (from O29).
      The three spatial dimensions are the representation-theoretic image of $\Im\mathbb{H}
  \cong \su(2)$, confirming the causal direction: $\Sigma_c(n_3) = 3 \Rightarrow
  \dim \Im\mathbb{H} = 3 \Rightarrow$ three spatial directions.
      Third, the bridge $\phi$ is unique up to positive scalar (Rigidity Lemma, Q7~Lemma~4.3),
      so the metric coefficients are uniquely determined.
      The no-cross-terms condition $k_X k_Z = k_Y k_Z = 0$ is proved analytically
      (Proposition~6.1).

      \textbf{Q9}~\cite{Beau2026q9} proves hypothesis~[H-lift] as a theorem by the method of
      generator suppression: the modulation generator is suppressed at rate $O(q^{-1/2})$ on the
      admissible sector, and coercivity of the limit form forces $A_H > 0$ independently of
      bridge existence.
      This makes Q5b Theorems~5.2 and~6.1 unconditional with respect to [H-lift].

      \textbf{Q8}~\cite{BeauQ8} establishes $\AZ = C_{\su(2)} = 2$ for the central (vertical)
      direction via Casimir rigidity.
      The Heisenberg commutator $[\tilde{X}, \tilde{Y}] = \tilde{Z}$ has a unique algebraic
      source; under the equivariant bridge $\phi$, its coefficient is constrained to be
      proportional to the unique $\su(2)$-invariant quadratic form on $\Sym^2(V_\rho)$, namely
      the Casimir $C_{\su(2)} = 2 \cdot \mathrm{Id}$.
      The normalisation is fixed with no free scalar.
      This closes Q5b open problem O2.

      \textbf{Q10}~\cite{Beau2026q10} proves $\AH \to 2$ as $q \to \infty$ under the O-series
      spectral universality hypothesis~[U].
      The argument runs: character-independence of the admissibility weights forces
      $\su(2)$-isotropy of the effective quadratic form on $H_{\mathrm{eff}} = \Sym^2(V_\rho)$;
      the unique $\su(2)$-invariant quadratic form on this spin-1 module has eigenvalue 2
      (Casimir).
      Numerical confirmation from the O25 campaign gives $|A_H(q) - 2| \approx 0.264 q^{-0.44}$
      for $q \in \{61, 101, 151\}$.

      \textbf{U1}~\cite{Beau2026u1} proves hypothesis~[U] with rate $\varepsilon(q) = O(q^{-1/2})$
      from the Weil-generator Lipschitz bound and the BFS--Carnot--Carath\'{e}odory convergence.
      This closes [U] as a theorem rather than a structural assumption.

      \textbf{Q11}~\cite{Beau2026q11} closes the sole remaining open coefficient $\Atau$.
      The cascade increment operator $T$ internalises no new irreducible representation; by
      Schur's lemma, the $\tau$-sector carries the trivial $\SU(2)$-representation and the
      temporal quadratic form must be $\SU(2)$-invariant.
      Since $T$ is the same cascade that fixes the spatial bridge normalisation (Q8), no
      independent normalisation arises: $\Atau = C_{\su(2)} = 2$.
      This gives the isotropic Lorentzian co-metric:
      \begin{equation}
        \label{eq:metric}
        \geff = 2\eta^{\mu\nu},
        \qquad \eta = \mathrm{diag}(-1, +1, +1, +1).
      \end{equation}
      The conformal factor of 2 is not a choice of units; it is the eigenvalue of the $\su(2)$-Casimir
      on the spin-1 module $\Sym^2(V_\rho)$, governing all four metric components uniformly.
      After absorbing $\Omega^2 = 2$ into the normalisation of the propagation speed, the
      effective metric is exactly Minkowski.

    \subsection{Closure papers: W1 and H2}
      \label{sec:closure}

      \textbf{W1}~\cite{Beau2026w1} proves hypothesis~[H-w] (convergence of the admissibility
      weights $a_q(s) \to A > 0$) as a structural consequence of the fibre-invariance of
      admissible observables, made quantitative by U1.
      After W1, the open Q5a hypotheses reduce to [H1], [H-E1], and~[C].

      \textbf{H2}~\cite{Beau2026h2} proves hypothesis~[H2] (strong convergence of the rescaled
      Weil generators $\hat{X}_q \to x$ and $\hat{P}_q \to -i\partial_x$) via a quantitative
      Poisson aliasing lemma and a discrete Sobolev identity.
      Together with Q5a-O2, H2 completes the closure of all Q5a hypotheses on the admissible
      sector.

      \begin{table}[ht]
        \centering
        \caption{Summary of constituent papers by internal phase.
        Status: P = proved; S = structural; C = conditional on Q5a hypotheses.}
        \label{tab:papers}
        \renewcommand{\arraystretch}{1.2}
        \begin{tabular}{@{}lllll@{}}
          \toprule
          Phase & Papers & Central output & Status \\
          \midrule
          Discrete-to-continuum & Q5a, Q5a-O2, H2 & $\Lpi = -A\partial_x^2$; [H-E1],[C],[H2] closed & P/S \\
          Dimensional promotion & Q5b & $\Dhom = 4$; Carnot convergence & S \\
          Metric extraction & Q5b + Q9 & $\geff \propto A_{\mu\nu}$; signature $(-, +, +, +)$ & P \\
          Coefficient determination & Q7, Q8, Q10, U1 & $\AZ = \AH = 2$ & P/S \\
          Metric closure & Q11, W1 & $\Atau = 2$; $\geff = 2\eta^{\mu\nu}$ & S/P \\
          Integrative output & Q6b & $\Pi_q \to \Leff \to \geff \to G_{\mu\nu}$ & C \\
          \bottomrule
        \end{tabular}
      \end{table}

  \section{Inputs from Upstream Results}
    \label{sec:inputs}

    The emergent geometry sub-programme takes the following inputs from Branch~I and from
    Presentation Note~1 (spectral admissibility).

    \begin{enumerate}
      \item \textbf{Weil representation on $L^2(\mathbb{Z}/q\mathbb{Z})$.}
      The identification of the admissible fibre as $V_\rho$ is a theorem of Foundation~M
      and HeisenbergStructure~\cite{Beau2026HeisenbergStructure,Beau2026m}.

      \item \textbf{Born--Infeld admissibility constraint.}
      The bounded-flux condition $A_n \leq c_{\mathrm{BI}}/\sqrt{\lambda_n}$ is the Q5a starting point.
      Its origin in the BI action is established in Branch~I.

      \item \textbf{Spin-$\tfrac{1}{2}$ sector and $\Sigma_c(n_3) = 3$.}
      From Presentation Note~1 (O23, O27, O29): the admissible sector is $V_\rho \cong \mathbb{C}^2$;
      the threshold dimension is $\Sigma_c(n_3) = 3$ by quaternionic minimality;
      $\Im\mathbb{H} \cong \su(2)$ is the three-dimensional spatial algebra.
      These are the inputs to Q7's identification of $H_{\mathrm{eff}} = \Sym^2(V_\rho)$ as
      the effective projection space.

      \item \textbf{O-series universality data.}
      The stabilisation of $\delta_{\mathrm{pair}}(q)$ across the O12--O25 campaign provides
      the numerical support for [H-E1] and for the spectral universality [U] required by Q10.
    \end{enumerate}

  \section{Outputs to Downstream Branches}
    \label{sec:outputs}

    Table~\ref{tab:outputs} lists the principal outputs of the sub-programme and their
    downstream consumers.

    \begin{table}[ht]
      \centering
      \caption{Principal outputs and downstream consumers.}
      \label{tab:outputs}
      \renewcommand{\arraystretch}{1.3}
      \begin{tabular}{@{}p{5.5cm}p{4.5cm}l@{}}
        \toprule
        Output & Consumer & Status \\
        \midrule
        $\geff = 2\eta^{\mu\nu}$ & Gravity, Q12, Q13 & S (Q11) \\
        Lorentzian signature $(-, +, +, +)$ & All Branch~III & P (Q5b + Q9) \\
        $\Leff$ on $\mathbb{R}_\tau \times \Heis_3(\mathbb{R})$ & Q6b, Gravity, Q12 & C \\
        $H_{\mathrm{eff}} = \Sym^2(V_\rho) \cong \mathbb{C}^3$ & Q7, Q8, Q14 & P \\
        Three spatial directions from $\Im\mathbb{H}$ & Q7, gauge notes & P \\
        Schwarzschild exterior metric & Q6b & C \\
        Einstein equations as consistency conditions & Q6b, Gravity & C \\
        $\ell_{\mathrm{sp}}$ (spectral length scale) & Gravity~\cite{Beau2026h} & S \\
        \bottomrule
      \end{tabular}
    \end{table}

  \section{Status of Results}
    \label{sec:status}

    \subsection{Proved results}
      \label{sec:proved}

      \begin{enumerate}
        \item \textit{$\Dhom = 4$} (Q5b Theorem~2.1): the Bass--Guivarc'h theorem for
        $\Heis_3(\mathbb{R})$; the homogeneous dimension is 4.

        \item \textit{Carnot convergence} (Q5b Theorem~3.2, following Pansu): the BFS shell
        stratification of $G_q$ converges to the Carnot--Carath\'{e}odory sphere foliation of
        $\Heis_3(\mathbb{R})$ in the pre-saturation regime.

        \item \textit{$[\text{H-lift}]$} (Q9 Theorem~1.2): the admissibility filter concentrates
        on the kinetic sector of $d\pi_1(\Delta_H)$ at rate $O(q^{-1/2})$, making $\Lpi$ the
        kinetic shadow of $\Delta_H$ without independent hypothesis.

        \item \textit{Metric extraction} (Q5b Theorem~5.2, unconditional via Q9): the effective
        co-metric $\geff \propto A_{\mu\nu}$ is extracted from the principal symbol of $\Leff$
        under the Q5a hypotheses.

        \item \textit{Lorentzian signature} (Q5b Theorem~6.1): signature $(-, +, +, +)$ is forced by
        the Born--Infeld admissibility constraint via exclusion of all other signature classes.

        \item \textit{Bridge rigidity} (Q7 Lemma~4.3): the bridge
        $\phi: \Sym^2(V_\rho) \xrightarrow{\sim} W_{\mathrm{sp}}$ is unique up to positive scalar.

        \item \textit{No-cross-terms} (Q7 Proposition~6.1): the cross terms $k_X k_Z = k_Y k_Z = 0$
        are proved analytically from $[\mathcal{F}_c, L_{\mathrm{Weil}}] = 0$.

        \item \textit{$\AZ = 2$} (Q8 Theorem~1.1): Casimir rigidity of the central direction;
        $\AZ = C_{\su(2)}|_{\Sym^2(V_\rho)} = 2 \cdot \mathrm{Id}$ with no free normalisation scalar.

        \item \textit{Spectral universality [U]} (U1): $\varepsilon(q) = O(q^{-1/2})$; proved from
        the Weil-generator Lipschitz bound.

        \item \textit{$\AH \to 2$} (Q10 Theorem~1.1): character-independence forces
        $\su(2)$-isotropy of the effective form; the Casimir gives value 2 in the limit.

        \item \textit{$[\text{H-w}]$} (W1): admissibility weight convergence $a_q(s) \to A > 0$
        proved as a structural consequence of fibre-invariance, quantified by U1.

        \item \textit{$[\text{H2}]$} (H2): strong convergence of rescaled Weil generators via
        Poisson aliasing lemma and discrete Sobolev identity.

        \item \textit{$[\text{H-E1}]$ and $[\text{C}]$ on the admissible sector} (Q5a-O2): proved
        via spectral atomicity of $\mathrm{ran}(\Pi_q)$.

        \item \textit{$\Atau = 2$} (Q11 Theorem~1.1): temporal Casimir rigidity; cascade
        internalization lemma, Schur rigidity, and Q8 normalisation close $\Atau$.

        \item \textit{Metric closure $\geff = 2\eta^{\mu\nu}$} (Q11 Corollary~6.1): combination of
        $\Atau = 2$ (Q11), $\AH = 2$ (Q10), $\AZ = 2$ (Q8), and Q5b Theorem~6.1.
      \end{enumerate}

    \subsection{Structural results}
      \label{sec:structural}

      \begin{enumerate}
        \item \textit{$H_{\mathrm{eff}} = \Sym^2(V_\rho)$} (Q7 Proposition~4.1, from O29):
        the effective projection space is the unique irreducible three-dimensional $\su(2)$-module.

        \item \textit{Three spatial dimensions from $\Im\mathbb{H}$} (Q7, O23): the three spatial
        directions are the image of $\Im\mathbb{H} \cong \su(2)$; three-dimensionality is not
        postulated.

        \item \textit{Schwarzschild exterior metric} (Q6b): flux conservation within the effective
        geometry forces the Schwarzschild metric as the unique stationary spherically symmetric
        exterior solution; conditional on Q5a hypotheses.

        \item \textit{Einstein equations as consistency conditions} (Q6b): the Einstein--Hilbert
        action is identified as the spectral entropy functional of the Gravity paper; conditional
        on Q5a hypotheses.

        \item \textit{$\AZ$ is not a free parameter} (Q8 §6.3): the value $\AZ = 2$ is fixed by
        the Casimir of the unique irreducible $\su(2)$-module of dimension 3; it is a structural
        identity, not a fitted constant.
      \end{enumerate}

    \subsection{Numerical evidence}
      \label{sec:numerical}

      \begin{enumerate}
        \item \textit{$|A_H(q) - 2| \approx 0.264 q^{-0.44}$} (Q7, Q10): confirmed from the O25
        campaign for $q \in \{61, 101, 151\}$; consistent with convergence to 2 with rate
        $O(q^{-1/2})$.

        \item \textit{Mosco tightness (Conjecture~[C])} (Q5a): spectral tail mass satisfies
        $\mathcal{E}_q(q/3)/\|f_q\|^2 < 2.1 \times 10^{-5}$ for $q \in \{29, 61, 101, 151\}$,
        with values decreasing toward machine precision.
        Closed analytically on the admissible sector by Q5a-O2.
      \end{enumerate}

    \subsection{Open hypotheses}
      \label{sec:open-hyp}

      \begin{enumerate}
        \item \textit{Hypothesis~[H1] on the full function space.}
        Q5a Hypothesis~[H1] (embedding $\iota_q: \mathcal{C}_q \hookrightarrow L^2(\mathbb{R})$
        with uniform control) is not needed in full generality on the admissible sector
        (Q5a-O2), but its extension to the full $L^2$ context remains open.
        This does not block any downstream result.

        \item \textit{Bridge existence at finite $q$} (Q8-O1).
        Q7 Conjecture~4.7 / Q8-O1: does an explicit $\su(2)$-equivariant isomorphism
        $\phi_q: \Sym^2(V_\rho) \xrightarrow{\sim} W_{\mathrm{sp}}$ exist at each prime $q$,
        not only in the asymptotic limit?
        The programme consequences ($\AZ = 2$, $\AH \to 2$, $\Atau = 2$) are established without
        it; this is a secondary structural open problem.
      \end{enumerate}

  \section{Open Deliverables}
    \label{sec:open}

    Two open items define the boundary of the sub-programme as of this note.

    \textbf{Open deliverable 1: Bridge existence at finite $q$.}
    The asymptotic results ($\AZ = \AH = \Atau = 2$) and the metric closure
    $\geff = 2\eta^{\mu\nu}$ are established in the $q \to \infty$ limit.
    The finite-$q$ equivariant bridge $\phi_q$ — required for a fully explicit
    representation-level identification at each prime — remains an open structural problem.
    Its resolution would strengthen the narrative but would not change any published result.

    \textbf{Open deliverable 2: [H1] on the full $L^2$ space.}
    Q5a-O2 closes [H1] on the admissible sector.
    The full-space version is not needed by any result in the programme and is listed here for
    completeness.

  \section{Conclusion}
    \label{sec:conclusion}

    The emergent geometry sub-programme derives the effective Lorentzian metric
    $\geff = 2\eta^{\mu\nu}$ from the admissibility filter $\Pi_q$ acting on the Weil
    representation of $\Heis_3(\mathbb{Z}/q\mathbb{Z})$.

    The chain has five stages, each resolved by distinct papers: the discrete-to-continuum
    Mosco limit (Q5a, Q5a-O2, W1, H2); the dimensional promotion via Carnot geometry (Q5b);
    the metric extraction and signature selection (Q5b + Q9); the determination of the metric
    coefficients $\AZ$, $\AH$, $\Atau$ via Casimir rigidity (Q7--Q11, U1); and the
    assembly into the effective spacetime statement (Q6b).
    All four metric coefficients equal 2, a single representation-theoretic datum: the
    eigenvalue of the $\su(2)$-Casimir on the spin-1 module $\Sym^2(V_\rho)$.

    The sub-programme closes the reconstruction problem: $\geff$ is not postulated but
    derived.
    The Lorentzian signature is not postulated but excluded by elimination.
    The value of the conformal factor is not fitted but forced by the Casimir.

    This note does not address the dynamics of $\geff$.
    The question of why the effective metric satisfies the Einstein equations is handled by the
    spectral gravity sub-programme (Presentation Note~4), which takes $\geff = 2\eta^{\mu\nu}$
    as its starting point.

    \newpage
    \bibliographystyle{plain}
    \bibliography{cosmochrony-bibliography}

\end{document}
